\documentclass[10pt,a4paper]{amsart}
\usepackage[margin=19mm]{geometry}
\usepackage{amsmath,amssymb,mathtools}
\usepackage[scr=rsfs,cal=euler]{mathalfa}
\usepackage{xfrac}
\usepackage[unicode,hidelinks,bookmarks=false]{hyperref}
\newcommand{\ie}{i.e.\ }
\newcommand{\cf}{cf.\ }
\newcommand{\resp}{respectively\ }
\newcommand{\Subsection}{\subsection}
\providecommand{\nobreakdash}{\nobreak\hbox{-}}
\setlength{\emergencystretch}{2em}
\multlinegap0pt
\usepackage{amssymb}
\usepackage[shortlabels]{enumitem}
\usepackage{float}
\usepackage{xcolor}
\newtheorem{assumption}{Assumption}
\newtheorem{theorem}{Theorem}
\usepackage{cleveref}
\crefname{assumption}{Assumption}{Assumptions}
\crefname{theorem}{Theorem}{Theorems}
\newcommand{\Embb}{\mathbb{E}}
\DeclareMathOperator{\Law}{Law}
\newcommand{\cF}{\mathcal{F}}
\newcommand{\cP}{\mathcal{P}}
\title{Global Multi-Maturity SPX--VIX Calibration Beyond Markovian Stitching}
\author{Atithi Acharya, Yue Sun, Brandon Augustino, Shouvanik Chakrabarti, Shree Hari Sureshbabu, Charlie Che}
\date{Source: arXiv:2609.04087v1 (CC BY 4.0)}
\begin{document}
\maketitle

\noindent\emph{Source and licence.} Global Multi-Maturity SPX--VIX Calibration Beyond Markovian Stitching. Version arXiv:2609.04087v1 (CC BY 4.0), distributed by arXiv under the Creative Commons Attribution 4.0 International licence, http://creativecommons.org/licenses/by/4.0/, as recorded in the arXiv licence metadata for this exact version and shown on the abstract page https://arxiv.org/abs/2609.04087v1. Full licence notice: \texttt{licenses/arxiv-2609.04087-CC-BY-4.0.txt}.\par\smallskip

\noindent\textit{Editorial scope.} Counted unit: the article's theorem thm:strict\_inclusion (source0.tex lines 354--368). The article's complete original proof of it is delivered with it (source0.tex lines 1114--1148, one \texttt{proof} environment of 35 lines, which the article places away from the statement). No mathematics is written, repaired or re-derived: the statement and the proof are the article's own lines, in the article's own notation, and no retained line is substituted or re-flowed.  Retained uncounted as the source-local context the proof needs: lines 166--179; lines 319--324; lines 326--345; lines 330--335.  Nothing was omitted from any retained span, so the package carries no omission record.  No retained line contains a citation, so the wrapper carries no bibliography and no bibliography entry.  No retained line contains a figure environment, an image-inclusion command or drawing code, so no image is delivered and the package declares no asset dependency.  Editorial formatting only, in addition: an AMS \texttt{amsart} preamble that restates the article's own declarations and macros used by the retained lines, namely the environments \texttt{assumption}, \texttt{theorem} on separate counters, together with the standard packages those lines need.  The retained environments print in this document as Assumption 1, Theorem 1 in order of appearance, whereas the article prints the same items with the numbering of its own full document.  The complete native source of the article as distributed in the arXiv e-print for this exact version is bundled unchanged as \texttt{sources/209342/source0.tex}.
\begin{assumption}\label{ass:marginals}
The given marginals satisfy
\begin{equation}
    \Embb^{\mu_{S_i}}[S_i] = S_0, \quad \Embb^{\mu_{S_i}}[|\ln S_i|] < \infty, \quad i = 1, \ldots, m;
\end{equation}
\begin{equation}
    \Embb^{\mu_{V_i}}[V_i^2] < \infty, \quad i = 1, \ldots, m-1.
\end{equation}
Moreover, the consistency condition
\begin{equation}\label{eq:consistency_global}
    \Embb^{\mu_{V_i}}[V_i^2] = \Embb^{\mu_{S_{i+1}}}[L(S_{i+1})] - \Embb^{\mu_{S_i}}[L(S_i)], \quad i = 1, \ldots, m-1
\end{equation}
holds, and $\mu_{S_i}$, $\mu_{S_{i+1}}$ are in convex order for each $i$.
\end{assumption}

\begin{equation}\label{eq:markovization}
\mathsf M\mu
:=
\nu_1^\mu(ds_1,dv_1,ds_2)
\prod_{i=2}^{m-1}\mu(dv_i,ds_{i+1}\mid s_i).
\end{equation}

\begin{theorem}[Feasibility equivalence and block-preserving Markovization]
\label{thm:feasibility_equivalence}
Assume regular conditional distributions exist and adjacent blocks use the same prescribed shared SPX marginal.
For every $\mu\in\cP_{\mathrm{full}}$,
\begin{equation}\label{eq:block_preservation}
    \mathsf M\mu\in\cP_{\mathrm{stitch}}\subseteq\cP_{\mathrm{full}},
    \qquad
    \Law_{\mathsf M\mu}(S_i,V_i,S_{i+1})
    =\Law_\mu(S_i,V_i,S_{i+1})
\end{equation}
for every $i$.
Consequently,
\begin{equation}\label{eq:feasibility_equivalence}
\cP_{\mathrm{full}}\neq\emptyset
\quad\Longleftrightarrow\quad
\cP_{\mathrm{stitch}}\neq\emptyset
\quad\Longleftrightarrow\quad
\cP_i\neq\emptyset\ \text{for every }i.
\end{equation}
\end{theorem}

\begin{equation}
    \mathsf M\mu\in\cP_{\mathrm{stitch}}\subseteq\cP_{\mathrm{full}},
    \qquad
    \Law_{\mathsf M\mu}(S_i,V_i,S_{i+1})
    =\Law_\mu(S_i,V_i,S_{i+1})
\end{equation}

\begin{theorem}[Strict law-class inclusion and local non-identification]
\label{thm:strict_inclusion}
For $m\geq3$, there exist marginals satisfying \cref{ass:marginals} for which
\begin{equation}
    \cP_{\mathrm{stitch}}\subsetneq\cP_{\mathrm{full}}.
\end{equation}
More precisely, there is a law $\mu\in\cP_{\mathrm{full}}$ with $\mu\neq\mathsf M\mu$ although all identities in~\eqref{eq:block_preservation} hold.
Thus every integrable block-local payoff $f_i(S_i,V_i,S_{i+1})$ has the same price under $\mu$ and $\mathsf M\mu$, while some bounded cross-period payoff $f$ has different prices under the two laws.
Moreover, $\mu=\mathsf M\mu$ if and only if
\begin{equation}\label{eq:stitched_ci}
    (V_i,S_{i+1})\ \perp\ \cF_{i-1}\mid S_i,
    \qquad i=2,\ldots,m-1,
\end{equation}
up to almost-sure equality of the corresponding kernels.
\end{theorem}

\begin{proof}[Proof of \cref{thm:strict_inclusion}]
The inclusion follows from \cref{thm:feasibility_equivalence}.
For strictness, take $m=3$, set $S_0=S_1=100$, and let $V_1$ be equally likely to equal $v_L=0.15$ or $v_H=0.55$.
For $(v,a)=(v_L,8)$ or $(v_H,30)$, set
\[
q(v,a):=\frac{v^2}{-(2/\tau)\log(1-(a/100)^2)}\in(0,1/2)
\]
and, conditional on $V_1=v$, assign probabilities $q(v,a)$, $1-2q(v,a)$, and $q(v,a)$ to $S_2=100-a$, $100$, and $100+a$.
Symmetry gives $\Embb[S_2\mid S_1,V_1]=S_1$, while
\[
q(v,a)\left[-\frac{2}{\tau}\log(1-a/100)-\frac{2}{\tau}\log(1+a/100)\right]=v^2,
\]
which gives the first dispersion identity.

Set $V_2=0.20$ when $V_1=v_L$ and $V_2=0.60$ when $V_1=v_H$ and, conditionally on the history, let
\[
S_3=S_2(1+\varepsilon\delta(V_2)),
\qquad
\delta(v):=\sqrt{1-e^{-v^2\tau}},
\]
where $\varepsilon$ is conditionally uniform on $\{-1,1\}$.
Then $\Embb[S_3\mid\cF_2]=S_2$ and
\[
\Embb[L(S_3/S_2)\mid\cF_2]
=-\frac{1}{\tau}\log(1-\delta(V_2)^2)=V_2^2,
\]
so the resulting finite-support law $\mu$ belongs to $\cP_{\mathrm{full}}$.

Both $V_1$ branches reach $S_2=100$ with positive probability, but $V_2$ identifies the branch there.
Hence $V_2\not\perp V_1\mid S_2$, so~\eqref{eq:stitched_ci} fails and $\mu\notin\cP_{\mathrm{stitch}}$.
Its Markovization belongs to $\cP_{\mathrm{stitch}}$ and preserves both adjacent triple laws by \cref{thm:feasibility_equivalence}.
Because the two finite laws differ, the indicator of any atom on which their masses differ is a bounded separating payoff.
Finally, equality $\mu=\mathsf M\mu$ is equivalent to equality of the full-history kernels and the $S_i$-conditional kernels in~\eqref{eq:markovization}, which is exactly~\eqref{eq:stitched_ci}.
For $m>3$, append degenerate feasible periods.
\end{proof}

\end{document}
